\documentclass[12pt,a4paper]{article}
\usepackage{../latex/dsdlab-report}

\begin{document}
	\reporttitlepage{۸}{\lr{ALU} اعداد مختلط}

	\section*{۸.۱ مقدمه و اهداف}
	این آزمایش از چند بخش تشکیل شده است. مطابق دستورکار، ابتدا باید پیمانه‌های گوناگون پیاده‌سازی و آزموده شوند و سپس با کنار هم گذاشتن آن‌ها یک رایانهٔ پایه ساخته شود که حافظه‌ای ۳۲ کلمه‌ای دارد. سه خواستهٔ صریح دستورکار عبارت‌اند از پیمانهٔ جمع و تفریق اعداد مختلط، پیمانهٔ ضرب اعداد مختلط، و واحدی که دستورها را از همان حافظه بخواند و عملیات مختلط را اجرا کند.

	افزون بر این‌ها، دستورکار یک قید مهم گذاشته است: «نبایستی از واحدهای جمع‌کننده یا ضرب‌کننده بیش از یک بار استفاده شود». این قید، شکل کلّی معماری طرح را تعیین می‌کند و در سراسر این گزارش به آن بازخواهیم گشت؛ چراکه رعایت آن به‌معنای پرهیز از هرگونه موازی‌سازی حسابی و تکیه بر اشتراک زمانی\LTRfootnote{Time-Multiplexed Resource Sharing} منابع است.

	\subsection*{روند انجام آزمایش}
	نخست دو بلوک حسابی مشترک، یعنی یک جمع‌کننده/تفریق‌کننده و یک ضرب‌کنندهٔ علامت‌دار، به‌صورت پیمانه‌های ترکیبی کوچک نوشته شدند. سپس چهار پیمانهٔ عملیاتی --- جمع و تفریق، ضرب، مزدوج و مربع اندازه --- هر یک به‌صورت یک ماشین حالت چندچرخه‌ای\LTRfootnote{Multi-Cycle Design} ساخته شدند که در هر چرخه تنها یک بار از بلوک حسابی مشترک استفاده می‌کند. برای هر پیمانه یک میز آزمون\LTRfootnote{Testbench} خودآزما نوشته و جداگانه اجرا شد. در گام بعد، همهٔ این عملیات در یک واحد حساب و منطق مختلط یکپارچه شدند و سرانجام این واحد به یک حافظهٔ دستور ۳۲ کلمه‌ای و یک واحد کنترل متّصل شد تا رایانهٔ خواسته‌شده کامل گردد. پس از تأیید همهٔ آزمون‌ها در \lr{ModelSim}، طرح در \lr{Quartus II} سنتز و تحلیل زمانی شد.

	\section*{۸.۲ نمایش اعداد و ساختار داده}
	\subsection*{الف) قالب عدد مختلط}
	هر عدد مختلط با دو بخش حقیقی و موهومی نمایش داده می‌شود و هر بخش یک عدد علامت‌دار هشت‌بیتی در نمایش متمّم دو\LTRfootnote{Two's Complement} است؛ یعنی هر بخش در بازهٔ \(-128\) تا \(+127\) قرار می‌گیرد. پهنای ورودی با پارامتر \lr{WIDTH} بیان شده و همهٔ پیمانه‌ها آن را می‌پذیرند، بنابراین طرح به عدد هشت وابستگی سخت ندارد.

	پهنای خروجی از رابطهٔ \(\text{\lr{RESULT\_WIDTH}} = 2\times\text{\lr{WIDTH}}+1\) به دست می‌آید که برای ورودی هشت‌بیتی برابر هفده بیت است. دلیل این انتخاب در ضرب مختلط نهفته است: حاصل‌ضرب دو عدد هشت‌بیتی علامت‌دار شانزده بیت می‌شود و بخش حقیقی نتیجه از تفریق دو چنین حاصل‌ضربی به دست می‌آید، پس برای جلوگیری از سرریز\LTRfootnote{Overflow} به یک بیت اضافی نیاز است. با هفده بیت، هیچ‌یک از چهار عملیات طرح نمی‌تواند از پهنای خروجی بیرون بزند.

	\subsection*{ب) قالب دستور و حافظه}
	حافظهٔ دستور همان ۳۲ کلمهٔ خواسته‌شدهٔ دستورکار است. هر کلمه سی و پنج بیت دارد و از الحاق یک کد عملیات سه‌بیتی با چهار میدان هشت‌بیتی ساخته می‌شود:
	\[
	\underbrace{\text{\lr{opcode}}}_{3},\;
	\underbrace{a_{\mathrm{re}}}_{8},\;
	\underbrace{a_{\mathrm{im}}}_{8},\;
	\underbrace{b_{\mathrm{re}}}_{8},\;
	\underbrace{b_{\mathrm{im}}}_{8}
	\]
	یعنی هر دستور، هم عملیات و هم هر دو عملوند مختلط خود را در دل خود دارد. این انتخاب، طرح را از داشتن فایل ثبات یا حافظهٔ داده بی‌نیاز می‌کند و تمرکز آزمایش را روی خود واحد حساب و اشتراک منابع نگه می‌دارد.

	\begin{table}[H]
		\centering
		\small
		\begin{tabular}{|c|l|l|}
			\hline
			\textbf{کد} & \textbf{نام} & \textbf{عملیات} \\ \hline
			\(000\) & \lr{CADD} & جمع دو عدد مختلط \\ \hline
			\(001\) & \lr{CSUB} & تفریق دو عدد مختلط \\ \hline
			\(010\) & \lr{CMUL} & ضرب دو عدد مختلط \\ \hline
			\(011\) & \lr{CONJ} & مزدوج عملوند نخست \\ \hline
			\(100\) & \lr{MAG2} & مربع اندازه‌ی عملوند نخست \\ \hline
			\(101\) & \lr{NOP}  & بدون عملیات \\ \hline
		\end{tabular}
		\caption{مجموعه‌ی دستورهای رایانه‌ی اعداد مختلط}
	\end{table}

	\section*{۸.۳ منابع حسابی مشترک}
	کلّ سخت‌افزار حسابی طرح تنها از دو بلوک ترکیبی کوچک تشکیل شده است: یک جمع‌کننده/تفریق‌کننده که با یک بیت انتخاب میان جمع و تفریق سوئیچ می‌کند، و یک ضرب‌کنندهٔ علامت‌دار. همهٔ عملیات مختلط طرح، بدون استثنا، از همین دو بلوک استفاده می‌کنند.

	\vcode{code/shared_arith.v}{دو بلوک حسابی مشترک}

	اهمّیّت این ساختار در برابر قید دستورکار روشن است. اگر ضرب مختلط به شیوهٔ متعارف و موازی نوشته می‌شد، به چهار ضرب‌کننده و دو جمع‌کننده نیاز داشت، چون
	\[
	(a+jb)(c+jd) = (ac-bd) + j(ad+bc).
	\]
	در این طرح به‌جای آن، همان یک ضرب‌کننده در چهار چرخهٔ پیاپی برای \(ac\)، \(bd\)، \(ad\) و \(bc\) به کار می‌رود و سپس همان یک جمع‌کننده در دو چرخهٔ دیگر تفریق و جمع پایانی را انجام می‌دهد. بهای این کار، طولانی‌تر شدن زمان اجرای هر عملیات است؛ سودش، رعایت قید دستورکار و کاهش چشمگیر مساحت.

	\section*{۸.۴ پیمانه‌های عملیاتی}
	\subsection*{الف) جمع و تفریق مختلط}
	این پیمانه چهار حالت دارد: بی‌کار، محاسبهٔ بخش حقیقی، محاسبهٔ بخش موهومی، و پایان. در حالت بی‌کار هر دو عملوند در ثبات‌های درونی نگه داشته می‌شوند؛ سپس در دو چرخهٔ پیاپی، همان یک جمع‌کننده نخست روی بخش‌های حقیقی و سپس روی بخش‌های موهومی کار می‌کند.

	پیش از ورود به جمع‌کننده، هر عملوند هشت‌بیتی با بسط علامت\LTRfootnote{Sign Extension} به هفده بیت تبدیل می‌شود تا مقدارهای منفی در پهنای بزرگ‌تر درست بمانند. آشکارسازی سرریز اما روی همان هشت بیت اصلی انجام می‌گیرد، زیرا پرسش معنادار این است که آیا نتیجه در پهنای ورودی جا می‌شود یا نه. برای جمع، سرریز هنگامی رخ می‌دهد که دو عملوند هم‌علامت باشند و علامت نتیجه با آن‌ها فرق کند؛ برای تفریق، هنگامی که دو عملوند ناهم‌علامت باشند و علامت نتیجه با عملوند نخست نخواند. پرچم سرریز نهایی، یای منطقی سرریز دو بخش است.

	\vcode{code/complex_addsub_mc.v}{پیمانه‌ی جمع و تفریق اعداد مختلط}

	\subsection*{ب) ضرب مختلط}
	این پیمانه هشت حالت دارد که چهار تای نخست آن‌ها چهار حاصل‌ضرب جزئی را یکی‌یکی می‌سازند و در ثبات نگه می‌دارند، و دو حالت بعدی با همان یک جمع‌کننده، تفریق بخش حقیقی و جمع بخش موهومی را انجام می‌دهند. بلوک ترکیبی انتخاب ورودی، در هر حالت عملوندهای درست را به ضرب‌کننده و جمع‌کننده می‌رساند؛ این بلوک در عمل همان تسهیم‌کننده‌های ورودی مسیر دادهٔ مشترک است.

	\vcode{code/complex_multiplier_mc.v}{پیمانه‌ی ضرب اعداد مختلط با یک ضرب‌کننده‌ی مشترک}

	\subsection*{ج) مزدوج و مربع اندازه}
	دو پیمانهٔ دیگر نیز با همین سبک ساخته شدند تا مجموعهٔ عملیات کامل شود. پیمانهٔ مزدوج، بخش حقیقی را دست‌نخورده نگه می‌دارد و بخش موهومی را منفی می‌کند؛ نکتهٔ آن این است که منفی‌سازی نیز به‌جای یک نفی‌کنندهٔ جداگانه، با همان جمع‌کنندهٔ مشترک و به‌صورت \(0-b\) انجام می‌شود تا قید دستورکار حتّی در این عمل ساده هم شکسته نشود.

	\vcode{code/complex_conjugate.v}{پیمانه‌ی مزدوج مختلط}

	پیمانهٔ مربع اندازه، مقدار \(|z|^2 = a^2 + b^2\) را در سه چرخه محاسبه می‌کند: دو چرخه برای دو مربع با همان یک ضرب‌کننده، و یک چرخه برای جمع آن‌ها. چون نتیجه ذاتاً حقیقی است، بخش موهومی خروجی همواره صفر گذاشته می‌شود.

	\vcode{code/magnitude_squared_mc.v}{پیمانه‌ی مربع اندازه}

	\section*{۸.۵ واحد حساب و منطق یکپارچه}
	پس از آزمودن جداگانهٔ پیمانه‌ها، همهٔ عملیات در یک واحد حساب و منطق مختلط گرد آمدند. این واحد شانزده حالت دارد که هر حالت، یک گام از یکی از عملیات را نمایندگی می‌کند؛ برای نمونه چهار حالت برای چهار حاصل‌ضرب جزئی \lr{CMUL} و دو حالت برای دو مؤلّفهٔ \lr{CADD}.

	نکتهٔ کلیدی معماری در همین‌جاست: در سراسر این واحد تنها \emph{یک} نمونه از ضرب‌کننده و \emph{یک} نمونه از جمع‌کننده ساخته می‌شود و همهٔ شانزده حالت از همان دو نمونه استفاده می‌کنند. حالت جاری تعیین می‌کند که کدام عملوندها به ورودی این دو بلوک برسند. به این ترتیب، افزودن عملیات تازه به این واحد هزینهٔ حسابی اضافه‌ای ندارد و تنها چند حالت و چند مسیر تسهیم‌کننده می‌افزاید.

	واحد حساب چهار پرچم وضعیّت نیز تولید می‌کند: سرریز، صفر بودن نتیجه، منفی بودن بخش حقیقی و منفی بودن بخش موهومی. هر چهار پرچم در حالت پایانی و هم‌زمان با نتیجه به‌روز می‌شوند تا هیچ‌گاه پرچمی متعلّق به عملیات پیشین در کنار نتیجهٔ تازه دیده نشود.

	\vcode{code/complex_alu_mc.v}{واحد حساب و منطق مختلط با منابع مشترک}

	\section*{۸.۶ حافظه‌ی دستور و واحد کنترل}
	حافظهٔ دستور یک آرایهٔ ۳۲ کلمه‌ای است که محتوای آن در خود توصیف سخت‌افزار مقداردهی می‌شود. این انتخاب عمدی است: با آن، \lr{ModelSim} و \lr{Quartus} هر دو دقیقاً همان برنامه را می‌بینند و نیازی به پروندهٔ جداگانهٔ مقداردهی حافظه نیست. خانه‌های استفاده‌نشده با دستور بدون عملیات پر می‌شوند تا هیچ خانه‌ای نامعیّن نماند.

	\vcode{code/instruction_memory.v}{حافظه‌ی دستور ۳۲ کلمه‌ای و برنامه‌ی آزمون}

	واحد کنترل، هفت حالت دارد که به‌ترتیب عبارت‌اند از بی‌کار، واکشی\LTRfootnote{Fetch} دستور، رمزگشایی\LTRfootnote{Decode} کد عملیات، خواندن میدان‌های عملوند، صدور پالس شروع به واحد حساب، انتظار تا پایان عملیات، و بازنویسی\LTRfootnote{Write-Back} نتیجه و پرچم‌ها. شمارندهٔ برنامه پس از هر بازنویسی یک واحد جلو می‌رود و با رسیدن به پایان برنامه، پرچم اتمام بالا می‌رود.

	\vcode{code/complex_processor.v}{واحد کنترل و اتّصال حافظه به واحد حساب}

	\subsection*{نکته‌ای درباره‌ی خواسته‌ی «پایپ‌لاین»}
	دستورکار در بند سوم خود، از «یک واحد پایپ‌لاین» سخن گفته است که دستورها را از حافظه بخواند و عملیات را به‌صورت پایپ‌لاین انجام دهد. لازم است در اینجا صریح باشیم: طرح حاضر دارای \emph{مراحل} جدا از هم به سبک پایپ‌لاین است --- واکشی، رمزگشایی، خواندن عملوند، اجرا و بازنویسی --- امّا این مراحل \emph{هم‌پوشانی} ندارند؛ یعنی دستور بعدی تنها پس از پایان کامل دستور جاری واکشی می‌شود. بنابراین طرح از نوع چندچرخه‌ای با مراحل مشخّص است و نه پایپ‌لاین با اجرای هم‌پوشان.

	این انتخاب، پیامد مستقیم قید دیگر همان صفحهٔ دستورکار است. پایپ‌لاین واقعی زمانی معنا دارد که چند دستور هم‌زمان در مراحل مختلف باشند؛ امّا وقتی تنها یک ضرب‌کننده و یک جمع‌کننده در کلّ طرح مجاز است، دو دستور حسابی نمی‌توانند هم‌زمان در مرحلهٔ اجرا باشند، زیرا هر دو به همان یک ضرب‌کننده نیاز دارند. رسیدن به هم‌پوشانی واقعی یا نیازمند تکثیر واحدهای حسابی است --- که قید دستورکار آن را منع می‌کند --- یا نیازمند منطق توقّف و مدیریت مخاطرهٔ ساختاری\LTRfootnote{Structural Hazard} است که در عمل، بهرهٔ پایپ‌لاین را برای این مجموعهٔ دستور از میان می‌برد. میان این دو خواسته، قید صریح و کمّی منابع بر توصیف کیفی پایپ‌لاین ترجیح داده شد.

	\section*{۸.۷ میزهای آزمون}
	برای هر پیمانه یک میز آزمون خودآزمای جداگانه نوشته شد و در پایان، یک میز آزمون یکپارچه کلّ رایانه را می‌آزماید. جدول زیر پوشش هر یک را خلاصه می‌کند.

	\begin{table}[H]
		\centering
		\footnotesize
		\begin{tabular}{|l|c|l|}
			\hline
			\textbf{میز آزمون} & \textbf{شمار وارسی} & \textbf{آنچه سنجیده می‌شود} \\ \hline
			جمع و تفریق & \(4\) & جمع مثبت، جمع منفی، تفریق، و حالت سرریز \\ \hline
			ضرب        & \(3\) & ضرب با نتیجه‌ی منفی، ضرب بزرگ، ضرب در عدد حقیقی \\ \hline
			مزدوج      & \(2\) & مزدوج با بخش موهومی مثبت و با بخش موهومی منفی \\ \hline
			مربع اندازه & \(3\) & مقدار متعارف، مقدار بزرگ، و ورودی صفر \\ \hline
			واحد حساب   & \(8\) & هر پنج عملیات به‌همراه پرچم‌های صفر و سرریز \\ \hline
			رایانه‌ی کامل & \(8\) & اجرای کامل برنامه‌ی حافظه و وارسی نتیجه‌ی هر دستور \\ \hline
		\end{tabular}
		\caption{پوشش شش میز آزمون آزمایش}
	\end{table}

	میز آزمون رایانهٔ کامل شیوهٔ متفاوتی دارد که ارزش توضیح دارد. این میز به‌جای انتظار برای زمان‌های از پیش تعیین‌شده، در هر لبهٔ پایین‌رونده پرچم اعتبار را می‌پاید و هرگاه نتیجه‌ای اعلام شود، با استفاده از شمارندهٔ برنامهٔ همراه آن نتیجه تشخیص می‌دهد که کدام دستور تمام شده و آن را با مقدار انتظاری خود می‌سنجد. این کار باعث می‌شود آزمون به طول متغیّر عملیات‌ها وابسته نباشد --- که در این طرح واقعاً متغیّر است، چون \lr{CADD} کوتاه و \lr{CMUL} بلند است --- و در پایان نیز بررسی می‌کند که دقیقاً هشت نتیجه دیده شده باشد.

	\vcode{code/tb_complex_processor.v}{میز آزمون یکپارچه‌ی رایانه‌ی اعداد مختلط}

	\section*{۸.۸ اجرا در \lr{ModelSim} و تحلیل کامل نتایج}
	هر آزمون با پروندهٔ فرمان جداگانهٔ خود اجرا می‌شود تا شکل موج هر پیمانه مستقل بررسی و ثبت گردد. فهرست کامل این فرمان‌ها در پرونده‌ای در کنار پروژه آمده و ترتیب اجرای آن‌ها چنین است: جمع و تفریق، ضرب، مزدوج، مربع اندازه، واحد حساب، و سرانجام رایانهٔ کامل.

	\docode{code/run_all.do}{پرونده‌ی فرمان اجرای شبیه‌سازی رایانه‌ی کامل}

	\subsection*{الف) کلید خواندن شکل موج‌ها}
	همهٔ شکل موج‌های این آزمایش با یک قاعدهٔ ساده خوانده می‌شوند. ساعت میز آزمون‌ها دورهٔ ۱۰ نانوثانیه دارد و هر پیمانه در هر لبهٔ فعّال ساعت دقیقاً یک گام از ماشین حالت خود را برمی‌دارد. بنابراین فاصلهٔ افقی میان پالس شروع و پالس اعتبار، شمار گام‌های آن عملیات را مستقیماً نشان می‌دهد.

	قرارداد سه سیگنالی همهٔ پیمانه‌ها یکسان است: با یک شدن \lr{start}، عملوندها در ثبات‌های درونی قفل می‌شوند و \lr{busy} بالا می‌رود؛ پیمانه گام‌های خود را یکی‌یکی پیش می‌برد؛ و در پایان، \lr{valid} برای دقیقاً یک پالس ساعت بالا می‌آید و \lr{busy} پایین می‌رود. مقدار خروجی تنها در همین لحظه معتبر است.

	چون هر گام یک چرخه است، شمار چرخه‌های هر عملیات را می‌توان پیش از هر شبیه‌سازی و تنها از روی فهرست حالت‌ها پیش‌بینی کرد. جدول زیر این پیش‌بینی را در کنار مقدار اندازه‌گیری‌شده از شکل موج می‌گذارد.

	\begin{table}[H]
		\centering
		\footnotesize
		\begin{tabular}{|l|l|c|c|}
			\hline
			\textbf{عملیات} & \textbf{گام‌های واحد حساب} & \textbf{چرخه} & \textbf{اندازه‌گیری} \\ \hline
			\lr{CADD} & جمع بخش حقیقی، جمع بخش موهومی، پایان & \(3\) & \(30\) نانوثانیه \\ \hline
			\lr{CSUB} & تفریق بخش حقیقی، تفریق بخش موهومی، پایان & \(3\) & \(30\) نانوثانیه \\ \hline
			\lr{CMUL} & چهار حاصل‌ضرب، تفریق، جمع، پایان & \(7\) & \(70\) نانوثانیه \\ \hline
			\lr{CONJ} & منفی‌سازی، پایان & \(2\) & \(20\) نانوثانیه \\ \hline
			\lr{MAG2} & دو مربع، جمع، پایان & \(4\) & \(40\) نانوثانیه \\ \hline
		\end{tabular}
		\caption{شمار چرخه‌ی هر عملیات؛ پیش‌بینی از روی حالت‌ها و اندازه‌گیری از شکل موج}
	\end{table}

	دو ستون آخر این جدول کاملاً بر هم منطبق‌اند و همین، بهای اشتراک منابع را عددی می‌کند: ضرب مختلط هفت چرخه می‌گیرد، در حالی که اگر چهار ضرب‌کنندهٔ موازی مجاز بود در یک چرخه انجام می‌شد. تفاوت هفت‌برابری زمان، بهایی است که برای حذف سه ضرب‌کننده پرداخته شده است.

	\subsection*{ب) پیمانه‌ی جمع و تفریق}
	\shot{screenshots/wave_addsub.png}{شکل موج پیمانه‌ی جمع و تفریق مختلط}
	\shotw{screenshots/transcript_addsub_1.png}{0.8}{رونوشت آزمون جمع و تفریق}
	\shotw{screenshots/transcript_addsub_2.png}{0.8}{ادامه‌ی رونوشت و گزارش قبولی چهار وارسی}

	در شکل موج چهار بازهٔ فعّالیّت پشت سر هم دیده می‌شود که هر یک ۳۵ نانوثانیه از پالس شروع تا پالس اعتبار طول می‌کشد. جدول زیر هر چهار وارسی را با بازهٔ زمانی، محاسبهٔ ریاضی و نتیجه می‌آورد.

	\begin{table}[H]
		\centering
		\footnotesize
		\begin{tabular}{|c|c|l|l|c|}
			\hline
			\textbf{\#} & \textbf{بازه (ns)} & \textbf{عملیات} & \textbf{محاسبه} & \textbf{سرریز} \\ \hline
			\(1\) & \(30\)--\(65\)   & \((3+j4)+(5+j2)\)     & \(3+5=8\)، \(4+2=6\)      & \(0\) \\ \hline
			\(2\) & \(80\)--\(115\)  & \((-3-j4)+(-5-j2)\)   & \(-3-5=-8\)، \(-4-2=-6\)  & \(0\) \\ \hline
			\(3\) & \(130\)--\(165\) & \((2+j3)-(7+j8)\)     & \(2-7=-5\)، \(3-8=-5\)    & \(0\) \\ \hline
			\(4\) & \(180\)--\(215\) & \((127+j0)+(1+j0)\)   & \(127+1=128\)             & \(1\) \\ \hline
		\end{tabular}
		\caption{چهار وارسی پیمانه‌ی جمع و تفریق با بازه‌ی زمانی و محاسبه‌ی هر یک}
	\end{table}

	درون هر بازه، الگوی سه‌گامی ماشین دیده می‌شود: در چرخهٔ نخست بخش حقیقی نتیجه در ثبات نوشته می‌شود، در چرخهٔ دوم بخش موهومی، و در چرخهٔ سوم پرچم اعتبار بالا می‌آید. همین ترتیب است که در نما به‌صورت دو پلّهٔ پیاپی روی خروجی‌ها و سپس یک پالس باریک اعتبار دیده می‌شود.

	وارسی چهارم بیش از همه ارزش توضیح دارد، چون در نگاه نخست متناقض به نظر می‌رسد: خروجی مقدار \(128\) را نشان می‌دهد و هم‌زمان پرچم سرریز بالاست. هر دو درست‌اند و به دو پرسش متفاوت پاسخ می‌دهند. خروجی هفده‌بیتی است و مقدار \(128\) به‌راحتی در آن جا می‌شود، پس عدد نمایش‌داده‌شده صحیح است. امّا منطق سرریز روی همان هشت بیت ورودی کار می‌کند و می‌پرسد آیا نتیجه در \emph{پهنای ورودی} جا می‌شود یا نه. هر دو عملوند مثبت‌اند و اگر جمع در هشت بیت انجام می‌شد به \(10000000_2\) می‌رسید که در متمّم دو برابر \(-128\) و منفی است. چون علامت نتیجه با علامت عملوندها نمی‌خواند، شرط سرریز جمع برقرار می‌شود:
	\[
	\text{سرریز} = \overline{(a_7 \oplus b_7)} \cdot (r_7 \oplus a_7).
	\]
	پس پیام واقعی این وارسی چنین است: «نتیجه ۱۲۸ است، ولی اگر می‌خواستید آن را دوباره در هشت بیت بریزید، معتبر نمی‌ماند».

	\subsection*{ج) پیمانه‌ی ضرب}
	\shot{screenshots/wave_multiplier.png}{شکل موج پیمانه‌ی ضرب مختلط}
	\shotw{screenshots/transcript_multiplier_1.png}{0.8}{رونوشت آزمون ضرب مختلط}
	\shotw{screenshots/transcript_multiplier_2.png}{0.78}{ادامه‌ی رونوشت و گزارش قبولی سه وارسی}

	اینجا اثر اشتراک منابع را می‌توان مستقیماً اندازه گرفت: هر بازه از پالس شروع تا پالس اعتبار ۷۵ نانوثانیه طول می‌کشد، یعنی هفت چرخهٔ کاری. جدول زیر نشان می‌دهد در هر چرخه چه چیزی محاسبه می‌شود.

	\begin{table}[H]
		\centering
		\footnotesize
		\begin{tabular}{|c|l|l|}
			\hline
			\textbf{چرخه} & \textbf{گام} & \textbf{کاری که انجام می‌شود} \\ \hline
			\(1\) & \(a\cdot c\) & ضرب‌کنندهٔ مشترک، بخش‌های حقیقی \\ \hline
			\(2\) & \(b\cdot d\) & ضرب‌کنندهٔ مشترک، بخش‌های موهومی \\ \hline
			\(3\) & \(a\cdot d\) & ضرب‌کنندهٔ مشترک، ضربدری \\ \hline
			\(4\) & \(b\cdot c\) & ضرب‌کنندهٔ مشترک، ضربدری \\ \hline
			\(5\) & \(ac-bd\) & جمع‌کنندهٔ مشترک در حالت تفریق \\ \hline
			\(6\) & \(ad+bc\) & جمع‌کنندهٔ مشترک در حالت جمع \\ \hline
			\(7\) & پایان & نوشتن نتیجه و بالا بردن پرچم اعتبار \\ \hline
		\end{tabular}
		\caption{هفت چرخه‌ی ضرب مختلط و کار هر چرخه}
	\end{table}

	سه وارسی این آزمون به‌ترتیب در بازه‌های \(30\)--\(105\)، \(120\)--\(195\) و \(210\)--\(285\) نانوثانیه انجام می‌شوند و جدول بعد هر یک را با تمام حاصل‌ضرب‌های جزئی می‌آورد تا روشن شود هر عدد خروجی از کجا آمده است.

	\begin{table}[H]
		\centering
		\footnotesize
		\begin{tabular}{|c|l|l|l|}
			\hline
			\textbf{\#} & \textbf{عملوندها} & \textbf{حاصل‌ضرب‌های جزئی} & \textbf{نتیجه} \\ \hline
			\(1\) & \((2+j3)\), \((-4+j5)\)  & \(ac=-8,\ bd=15,\ ad=10,\ bc=-12\)  & \(-8-15=-23\)، \(10-12=-2\) \\ \hline
			\(2\) & \((-10+j5)\), \((3-j8)\) & \(ac=-30,\ bd=-40,\ ad=80,\ bc=15\) & \(-30+40=10\)، \(80+15=95\) \\ \hline
			\(3\) & \((0+j2)\), \((0+j2)\)   & \(ac=0,\ bd=4,\ ad=0,\ bc=0\)       & \(0-4=-4\)، \(0+0=0\) \\ \hline
		\end{tabular}
		\caption{سه وارسی ضرب مختلط به‌همراه چهار حاصل‌ضرب جزئی هر یک؛ بازه‌های زمانی در متن آمده‌اند}
	\end{table}

	وارسی سوم عمداً انتخاب شده است، چون همان اتّحاد آشنای \(j\cdot j=-1\) را می‌آزماید: ضرب \(j2\) در \(j2\) باید \(-4\) بدهد و بخش موهومی آن صفر باشد. این تنها حالتی است که در آن بخش حقیقی نتیجه منفی می‌شود بدون آنکه هیچ‌یک از عملوندها بخش حقیقی داشته باشد، و درستی علامت‌گذاری تفریق \(ac-bd\) را می‌سنجد.

	\subsection*{د) پیمانه‌های مزدوج و مربع اندازه}
	\shot{screenshots/wave_conjugate.png}{شکل موج پیمانه‌ی مزدوج}
	\shotw{screenshots/transcript_conjugate_1.png}{0.72}{رونوشت آزمون مزدوج}
	\shotw{screenshots/transcript_conjugate_2.png}{0.78}{گزارش قبولی دو وارسی مزدوج}

	مزدوج کوتاه‌ترین عملیات طرح است و تنها دو چرخه می‌گیرد؛ در شکل موج، فاصلهٔ شروع تا اعتبار ۲۵ نانوثانیه است. دو وارسی آن، هر دو جهت تغییر علامت را می‌آزمایند: بازهٔ \(30\)--\(55\) ورودی \((-6+j7)\) را به \((-6-j7)\) و بازهٔ \(70\)--\(95\) ورودی \((4-j9)\) را به \((4+j9)\) تبدیل می‌کند. در هر دو حالت بخش حقیقی دست‌نخورده می‌ماند و تنها بخش موهومی از راه عمل \(0-b\) روی جمع‌کنندهٔ مشترک منفی می‌شود.

	\shot{screenshots/wave_mag2.png}{شکل موج پیمانه‌ی مربع اندازه}
	\shotw{screenshots/transcript_mag2_1.png}{0.7}{رونوشت آزمون مربع اندازه}
	\shotw{screenshots/transcript_mag2_2.png}{0.8}{گزارش قبولی سه وارسی مربع اندازه}

	\begin{table}[H]
		\centering
		\footnotesize
		\begin{tabular}{|c|c|l|l|}
			\hline
			\textbf{\#} & \textbf{بازه (ns)} & \textbf{ورودی} & \textbf{محاسبه} \\ \hline
			\(1\) & \(30\)--\(75\)   & \(3-j4\)   & \(3^2+(-4)^2=9+16=25\) \\ \hline
			\(2\) & \(90\)--\(135\)  & \(-10+j5\) & \((-10)^2+5^2=100+25=125\) \\ \hline
			\(3\) & \(150\)--\(195\) & \(0+j0\)   & \(0+0=0\) \\ \hline
		\end{tabular}
		\caption{سه وارسی مربع اندازه با بازه‌ی زمانی و محاسبه‌ی هر یک}
	\end{table}

	هر بازه ۴۵ نانوثانیه است که با چهار چرخهٔ این عملیات می‌خواند. نکتهٔ خواندنی دو وارسی نخست آن است که با وجود منفی بودن یکی از مؤلّفه‌های ورودی، نتیجه مثبت درمی‌آید؛ زیرا هر مؤلّفه پیش از جمع در خودش ضرب می‌شود و مربع هر عدد علامت‌دار نامنفی است. وارسی سوم نیز حالت مرزی ورودی صفر را می‌آزماید که تنها ورودی ممکن با نتیجهٔ صفر است.

	\subsection*{ه) واحد حساب یکپارچه}
	\shot{screenshots/wave_alu.png}{شکل موج واحد حساب و منطق مختلط}
	\shotw{screenshots/transcript_alu_1.png}{0.72}{رونوشت آزمون واحد حساب}
	\shotw{screenshots/transcript_alu_2.png}{0.78}{ادامه‌ی رونوشت واحد حساب}
	\shotw{screenshots/transcript_alu_3.png}{0.78}{گزارش قبولی هشت وارسی واحد حساب}

	هشت وارسی این آزمون، هر پنج عملیات را به‌همراه هر چهار پرچم می‌پوشانند. جدول زیر همهٔ آن‌ها را با بازهٔ زمانی و مقدار انتظاری هر پرچم گرد می‌آورد.

	\begin{table}[H]
		\centering
		\footnotesize
		\begin{tabular}{|c|c|l|l|c|c|c|c|}
			\hline
			\textbf{\#} & \textbf{بازه (ns)} & \textbf{عملیات} & \textbf{نتیجه} & \textbf{سرریز} & \textbf{صفر} & \textbf{حقیقی منفی} & \textbf{موهومی منفی} \\ \hline
			\(1\) & \(30\)--\(65\)   & \lr{CADD} \((3+j4)+(5+j2)\)   & \(8+j6\)   & \(0\) & \(0\) & \(0\) & \(0\) \\ \hline
			\(2\) & \(80\)--\(115\)  & \lr{CADD} منفی‌ها             & \(-8-j6\)  & \(0\) & \(0\) & \(1\) & \(1\) \\ \hline
			\(3\) & \(130\)--\(165\) & \lr{CSUB} \((2+j3)-(7+j8)\)   & \(-5-j5\)  & \(0\) & \(0\) & \(1\) & \(1\) \\ \hline
			\(4\) & \(180\)--\(255\) & \lr{CMUL} \((2+j3)(-4+j5)\)   & \(-23-j2\) & \(0\) & \(0\) & \(1\) & \(1\) \\ \hline
			\(5\) & \(270\)--\(295\) & \lr{CONJ} \((-6+j7)\)         & \(-6-j7\)  & \(0\) & \(0\) & \(1\) & \(1\) \\ \hline
			\(6\) & \(310\)--\(355\) & \lr{MAG2} \((3-j4)\)          & \(25\)     & \(0\) & \(0\) & \(0\) & \(0\) \\ \hline
			\(7\) & \(370\)--\(405\) & \lr{CADD} \((5-j5)+(-5+j5)\)  & \(0\)      & \(0\) & \(1\) & \(0\) & \(0\) \\ \hline
			\(8\) & \(420\)--\(455\) & \lr{CADD} \((127+j0)+(1+j0)\) & \(128\)    & \(1\) & \(0\) & \(0\) & \(0\) \\ \hline
		\end{tabular}
		\caption{هشت وارسی واحد حساب با بازه‌ی زمانی و مقدار هر چهار پرچم}
	\end{table}

	طول بازه‌ها در این جدول، همان جدول چرخه‌های بند الف را تأیید می‌کند: وارسی‌های جمع و تفریق ۳۵ نانوثانیه، ضرب ۷۵ نانوثانیه، مزدوج ۲۵ نانوثانیه و مربع اندازه ۴۵ نانوثانیه طول می‌کشند. یعنی همان پیمانه‌های جداگانه، پس از یکپارچه شدن در واحد حساب، دقیقاً همان شمار چرخه را نگه داشته‌اند.

	پرچم‌ها نیز الگوی روشنی دارند. پرچم صفر تنها در وارسی هفتم بالاست، چون تنها همان‌جاست که هر دو مؤلّفهٔ نتیجه صفر می‌شوند؛ توجّه شود که نتیجهٔ مربع اندازه در وارسی ششم بخش موهومی صفر دارد ولی بخش حقیقی آن ۲۵ است، پس پرچم صفر بالا نمی‌رود. پرچم سرریز تنها در وارسی هشتم بالاست و در وارسی چهارم که نتیجهٔ \(-23\) از تفریق دو حاصل‌ضرب به دست آمده پایین می‌ماند، زیرا منطق سرریز تنها برای عملیات جمع و تفریق هشت‌بیتی معنا دارد و پهنای هفده‌بیتی برای نتیجهٔ ضرب همیشه کافی است.

	\subsection*{و) رایانه‌ی کامل}
	\shot{screenshots/wave_processor.png}{شکل موج اجرای کامل برنامه‌ی حافظه}

	این شکل موج، گویاترین نمای کلّ آزمایش است و می‌توان کلّ برنامه را از روی آن بازخواند. ردیف شمارندهٔ نتیجه به‌ترتیب مقدارهای صفر تا هفت را می‌گیرد و در کنار هر یک، بخش حقیقی و موهومی نتیجه و وضعیّت پرچم‌ها دیده می‌شود. جدول زیر لحظهٔ دقیق پایان هر دستور و فاصلهٔ آن از دستور پیش را می‌آورد.

	\begin{table}[H]
		\centering
		\footnotesize
		\begin{tabular}{|c|l|c|c|c|}
			\hline
			\textbf{شمارنده} & \textbf{دستور} & \textbf{لحظه‌ی اعتبار (ns)} & \textbf{فاصله (ns)} & \textbf{چرخه} \\ \hline
			\(0\) & \lr{CADD} & \(135\) & ---     & \(9\) \\ \hline
			\(1\) & \lr{CADD} & \(225\) & \(90\)  & \(9\) \\ \hline
			\(2\) & \lr{CSUB} & \(315\) & \(90\)  & \(9\) \\ \hline
			\(3\) & \lr{CMUL} & \(445\) & \(130\) & \(13\) \\ \hline
			\(4\) & \lr{CONJ} & \(525\) & \(80\)  & \(8\) \\ \hline
			\(5\) & \lr{MAG2} & \(625\) & \(100\) & \(10\) \\ \hline
			\(6\) & \lr{CADD} & \(715\) & \(90\)  & \(9\) \\ \hline
			\(7\) & \lr{CADD} & \(805\) & \(90\)  & \(9\) \\ \hline
		\end{tabular}
		\caption{لحظه‌ی پایان هر دستور برنامه و فاصله‌ی آن از دستور پیشین}
	\end{table}

	ستون آخر این جدول با یک رابطهٔ ساده توضیح‌پذیر است. هزینهٔ هر دستور از دو بخش تشکیل می‌شود: سربار ثابت واحد کنترل و گام‌های متغیّر واحد حساب. با مقایسهٔ اعداد اندازه‌گیری‌شده با جدول بند الف به این رابطه می‌رسیم:
	\[
	\text{چرخه‌های هر دستور} = 6 + \text{گام‌های واحد حساب}.
	\]
	درستی آن را می‌توان برای هر چهار عملیات وارسی کرد: \(6+3=9\) برای جمع و تفریق، \(6+7=13\) برای ضرب، \(6+2=8\) برای مزدوج و \(6+4=10\) برای مربع اندازه. سربار ثابت شش‌چرخه‌ای نیز خود قابل تفکیک است: یک چرخه واکشی، یک چرخه رمزگشایی، یک چرخه خواندن عملوندها، یک چرخه صدور پالس شروع، یک چرخه برای دیدن پرچم اعتبار واحد حساب، و یک چرخه بازنویسی نتیجه.

	همین رابطه، ادّعای بند ۸.۶ دربارهٔ نبود هم‌پوشانی را نیز عددی می‌کند: اگر مراحل واقعاً هم‌پوشان بودند، فاصلهٔ میان دو دستور پیاپی به جای نه چرخه باید نزدیک به یک چرخه می‌شد. مشاهدهٔ فاصلهٔ ثابت نه‌چرخه‌ای میان دستورهای هم‌نوع، مستقیماً نشان می‌دهد که دستور بعدی تنها پس از پایان کامل دستور جاری آغاز می‌شود.

	جدول بعد، محتوای برنامه و محاسبهٔ هر نتیجه را می‌آورد تا روشن شود چرا اعداد روی شکل موج همان‌اند که هستند.

	\begin{table}[H]
		\centering
		\footnotesize
		\begin{tabular}{|c|l|l|l|c|}
			\hline
			\textbf{شمارنده} & \textbf{دستور و عملوندها} & \textbf{محاسبه} & \textbf{نتیجه} & \textbf{پرچم فعّال} \\ \hline
			\(0\) & \lr{CADD} \((3+j4)\), \((5+j2)\)   & \(3+5\)، \(4+2\)                     & \(8+j6\)    & --- \\ \hline
			\(1\) & \lr{CADD} \((-3-j4)\), \((-5-j2)\) & \(-3-5\)، \(-4-2\)                   & \(-8-j6\)   & هر دو منفی \\ \hline
			\(2\) & \lr{CSUB} \((2+j3)\), \((7+j8)\)   & \(2-7\)، \(3-8\)                     & \(-5-j5\)   & هر دو منفی \\ \hline
			\(3\) & \lr{CMUL} \((2+j3)\), \((-4+j5)\)  & \(-8-15\)، \(10-12\)                 & \(-23-j2\)  & هر دو منفی \\ \hline
			\(4\) & \lr{CONJ} \((-6+j7)\)              & تغییر علامت بخش موهومی               & \(-6-j7\)   & هر دو منفی \\ \hline
			\(5\) & \lr{MAG2} \((3-j4)\)               & \(9+16\)                             & \(25\)      & --- \\ \hline
			\(6\) & \lr{CADD} \((5-j5)\), \((-5+j5)\)  & \(5-5\)، \(-5+5\)                    & \(0\)       & صفر \\ \hline
			\(7\) & \lr{CADD} \((127+j0)\), \((1+j0)\) & \(127+1\)                            & \(128\)     & سرریز \\ \hline
		\end{tabular}
		\caption{برنامه‌ی هشت‌دستوری حافظه، محاسبه‌ی هر دستور و پرچم فعّال آن}
	\end{table}

	\shotw{screenshots/transcript_processor_1.png}{0.8}{آغاز رونوشت اجرای رایانه‌ی کامل}
	\shotw{screenshots/transcript_processor_2.png}{0.78}{وارسی نتیجه‌ی هشت دستور برنامه}
	\shotw{screenshots/transcript_processor_3.png}{0.75}{گزارش پایانی و قبولی هر هشت دستور}

	چند نکته در خواندن این نتایج اهمّیّت دارد. نخست آنکه دستور شمارهٔ چهار، یعنی مزدوج، تنها دستوری است که یک عملوند دارد؛ دو میدان دیگر دستور در حافظه صفر گذاشته شده‌اند و واحد حساب هم آن‌ها را نادیده می‌گیرد. دوم آنکه دستور شمارهٔ پنج نتیجه‌ای کاملاً حقیقی می‌دهد و بخش موهومی آن نه از محاسبه، بلکه با صفر گذاشتن صریح در حالت پایانی به دست می‌آید. سوم آنکه دو دستور پایانی عمداً برای آزمودن پرچم‌ها انتخاب شده‌اند: دستور ششم دو عدد قرینه را جمع می‌کند تا هر دو مؤلّفه صفر شوند و پرچم صفر بالا برود، و دستور هفتم جمعی می‌سازد که از بازهٔ هشت‌بیتی علامت‌دار بیرون می‌زند تا پرچم سرریز آزموده شود.

	سرانجام، شمارش کلّی نیز با اعداد جدول می‌خواند. مجموع چرخه‌های هشت دستور برابر
	\[
	9+9+9+13+8+10+9+9 = 76
	\]
	چرخه است و با احتساب چرخه‌های بازنشانی و پالس شروع در ابتدای اجرا، پرچم اتمام در لحظهٔ ۸۰۵ نانوثانیه بالا می‌رود و شبیه‌سازی در ۸۲۰ نانوثانیه پایان می‌یابد.

	\subsection*{ز) جمع‌بندی تحلیل}
	آنچه از بررسی شکل موج‌ها و رونوشت‌ها به دست می‌آید چنین است:
	\begin{itemize}
		\item شمار چرخهٔ هر عملیات، که پیش از شبیه‌سازی و تنها از روی فهرست حالت‌ها پیش‌بینی شده بود، در هر شش آزمون دقیقاً با اندازه‌گیری خوانده است.
		\item همهٔ بیست و هشت عدد خروجی، از حاصل‌ضرب‌های جزئی ضرب مختلط تا مربع اندازه، با محاسبهٔ دستی قابل بازسازی‌اند و هیچ‌یک نیازمند اعتماد کورکورانه به شبیه‌ساز نیست.
		\item چهار پرچم وضعیّت در حالت‌های مرزی آزموده شدند و رفتار به‌ظاهر متناقض وارسی سرریز، با تفکیک پهنای خروجی از پهنای ورودی توضیح داده شد.
		\item پیمانه‌های جداگانه پس از یکپارچه شدن در واحد حساب، همان شمار چرخه را نگه داشتند؛ یعنی یکپارچه‌سازی هزینهٔ زمانی تازه‌ای تحمیل نکرده است.
		\item فاصلهٔ ثابت میان دستورهای هم‌نوع در شکل موج رایانهٔ کامل، نبود هم‌پوشانی میان دستورها را به‌صورت عددی تأیید می‌کند.
	\end{itemize}

	\section*{۸.۹ سنتز و تحلیل زمانی}
	\subsection*{الف) مصرف منابع}
	طرح در \lr{Quartus II 13.0.1} و برای خانوادهٔ \lr{Cyclone IV GX} همگردانی شد. انتخاب قطعه روی حالت خودکار بود و ابزار قطعهٔ \lr{EP4CGX22CF19C6} را برگزید. نتیجه ۵۲۰ عنصر منطقی شامل ۴۹۷ تابع ترکیبی و ۲۹۸ ثبات اختصاصی، و ۴۹ پایه است.

	\shotw{screenshots/flow_summary.png}{0.8}{خلاصه‌ی منابع مصرفی رایانه‌ی اعداد مختلط}

	یک عدد این گزارش شایان توجّه ویژه است: شمار ضرب‌کننده‌های نه‌بیتی توکار\LTRfootnote{Embedded Multiplier} مصرف‌شده صفر است. یعنی سنتز‌کننده هیچ بلوک ضرب آمادهٔ تراشه را به کار نگرفته و تنها ضرب‌کنندهٔ طرح را با جدول‌های جست‌وجو ساخته است. این نکته از دو جهت اهمّیّت دارد: نخست آنکه تأیید می‌کند در کلّ طرح واقعاً بیش از یک ضرب‌کننده وجود ندارد، و دوم آنکه بخش بزرگی از ۴۹۷ تابع ترکیبی گزارش‌شده صرف همین یک ضرب‌کنندهٔ هشت در هشت شده است.

	شمار ۴۹ پایه نیز از واسط طرح قابل بازبینی است: ساعت، بازنشانی و شروع در ورودی، و دو نتیجهٔ هفده‌بیتی، شمارندهٔ پنج‌بیتی و شش پرچم در خروجی.

	\subsection*{ب) نمای سطح انتقال ثبات‌ها و ماشین حالت}
	نمای سطح انتقال ثبات‌ها ساختار مورد انتظار را نشان می‌دهد: بلوک حافظهٔ دستور در سمت چپ، ثبات‌های میان‌مرحله‌ای که دستور و عملوندهای رمزگشایی‌شده را نگه می‌دارند، بلوک واحد حساب در میانه، و ثبات‌های خروجی نتیجه و پرچم‌ها در سمت راست.

	\shot{screenshots/rtl_processor.png}{نمای سطح انتقال ثبات‌های رایانه‌ی اعداد مختلط}

	نمای ماشین حالت\LTRfootnote{State Machine Viewer} نیز هفت حالت واحد کنترل و گذارهای میان آن‌ها را نشان می‌دهد. حلقهٔ بازگشتی روی حالت انتظار، همان انتظار برای پایان عملیات واحد حساب است و حلقهٔ بزرگ‌تری که از بازنویسی به واکشی بازمی‌گردد، پیمایش دستورهای پیاپی برنامه را نشان می‌دهد.

	\shot{screenshots/state_machine_viewer.png}{نمای ماشین حالت واحد کنترل}

	\shot{screenshots/post_mapping.png}{نمای \lr{Technology Map Viewer (Post-Mapping)}}
	\shot{screenshots/post_fitting.png}{نمای \lr{Technology Map Viewer (Post-Fitting)}}

	\subsection*{ج) تحلیل زمانی}
	برخلاف چند آزمایش پیشین، برای این طرح یک پروندهٔ قیود طراحی\LTRfootnote{Synopsys Design Constraints} نوشته شد که ساعتی با دورهٔ ۲۰ نانوثانیه، معادل ۵۰ مگاهرتز، تعریف می‌کند و عدم‌قطعیّت ساعت را نیز استخراج می‌نماید.

	\docode{code/complex_processor.sdc}{پرونده‌ی قیود طراحی}

	تحلیل‌گر زمانی\LTRfootnote{TimeQuest Timing Analyzer} بیشینه‌بسامد\LTRfootnote{Maximum Frequency (Fmax)} را \(88.25\,\mathrm{MHz}\) گزارش می‌کند و این بار بیشینه‌بسامد مقیّد نیز همین مقدار است؛ یعنی محدودکنندهٔ سرعت، منطق درونی طرح است و نه نرخ تغییر پایه‌ها.

	\shotw{screenshots/fmax_summary.png}{0.85}{گزارش بیشینه‌بسامد رایانه‌ی اعداد مختلط}

	از این عدد می‌توان حاشیهٔ زمانی را بازسازی کرد: بلندترین مسیر داده برابر \(1/88.25\,\mathrm{MHz}\approx11.33\) نانوثانیه است، پس در برابر قید ۲۰ نانوثانیه‌ای حاشیه‌ای نزدیک به \(8.7\) نانوثانیه باقی می‌ماند و طرح قید خود را با اطمینان برآورده می‌کند. مجموع حاشیه‌های منفی در همهٔ تحلیل‌ها صفر است و حاشیهٔ نگهداری نیز مثبت می‌ماند.

	بیشینه‌بسامد این طرح به‌روشنی کمتر از آزمایش‌های پیشین است و دلیل آن هم آشکار است: مسیر بحرانی از ضرب‌کنندهٔ هشت در هشتی می‌گذرد که با جدول‌های جست‌وجو ساخته شده و پس از آن به جمع‌کنندهٔ هفده‌بیتی می‌رسد. این همان بهایی است که برای اشتراک منابع و پرهیز از ضرب‌کننده‌های موازی پرداخته می‌شود.

	با توجّه به برگزاری مجازی آزمایشگاه، ارزیابی این آزمایش بر پایهٔ شبیه‌سازی، سنتز و تحلیل زمانی نرم‌افزاری انجام شده و در پرونده‌های پیکربندی پروژه تخصیص پایه‌ای برای برد ثبت نشده است؛ ازاین‌رو ادّعایی دربارهٔ آزمون فیزیکی روی برد مطرح نمی‌شود.

	\section*{۸.۱۰ نتایج}
	\begin{itemize}
		\item پیمانه‌های جمع و تفریق و ضرب اعداد مختلط، مطابق خواستهٔ دستورکار، جداگانه پیاده‌سازی و هر یک با میز آزمون خودآزمای خود آزموده شدند.
		\item دو پیمانهٔ مزدوج و مربع اندازه نیز افزوده شدند تا مجموعهٔ عملیات واحد حساب کامل شود.
		\item قید دستورکار دربارهٔ استفادهٔ یک‌بارهٔ واحدهای حسابی به‌طور کامل رعایت شد: در کلّ طرح تنها یک ضرب‌کننده و یک جمع‌کننده وجود دارد و همهٔ عملیات با اشتراک زمانی از آن‌ها استفاده می‌کنند.
		\item صفر بودن ضرب‌کننده‌های توکار در گزارش سنتز، این ادّعا را به‌صورت مستقل تأیید می‌کند.
		\item رایانهٔ خواسته‌شده با حافظهٔ ۳۲ کلمه‌ای ساخته شد و برنامهٔ هشت‌دستوری آن، شامل هر پنج عملیات و دو حالت پرچمی، بدون خطا اجرا شد.
		\item در مجموع بیست و هشت وارسی در شش میز آزمون انجام شد و همگی با موفّقیّت گذشتند.
		\item سنتز با ۵۲۰ عنصر منطقی، ۲۹۸ ثبات و ۴۹ پایه انجام شد و تحلیل زمانی با قید ۵۰ مگاهرتزی بیشینه‌بسامد \(88.25\,\mathrm{MHz}\) و حاشیهٔ مثبت داد.
	\end{itemize}

	\section*{نتیجه‌گیری}
	در این آزمایش یک رایانهٔ پایه برای اعداد مختلط از پایین به بالا ساخته شد: نخست دو بلوک حسابی مشترک، سپس چهار پیمانهٔ عملیاتی، آنگاه یک واحد حساب یکپارچه و سرانجام واحد کنترلی که برنامه‌ای را از حافظه می‌خواند و اجرا می‌کند. آنچه شکل این معماری را تعیین کرد، قید صریح دستورکار دربارهٔ استفادهٔ یک‌بارهٔ واحدهای حسابی بود؛ قیدی که طرح را ناگزیر از چندچرخه‌ای شدن کرد و در ازای زمان اجرای بیشتر، مساحت را چشمگیر کاهش داد. همین قید بود که ما را واداشت میان دو خواستهٔ ناسازگار دستورکار --- اجرای پایپ‌لاین و اشتراک کامل منابع --- آگاهانه انتخاب کنیم و دلیل انتخاب خود را در گزارش بیاوریم. سرانجام، گزارش سنتز با صفر اعلام کردن ضرب‌کننده‌های توکار، نشان داد که ادّعای اشتراک منابع تنها ادّعایی در سطح کد نیست و در سخت‌افزار ساخته‌شده نیز برقرار است.
\end{document}
